%% ===================================================================
\section{Conclusion}
\label{sec:conclusion}
%% ===================================================================

\AEGIS provides a principled pre-deployment diagnostic for GNNs on safety-critical graphs: from a single $S_c$ computation, practitioners obtain the first-order maximally sensitive perturbation direction, per-edge vulnerability rankings, and per-node first-order sensitivity thresholds. The $S_c$ computation applies to any GNN with continuous edge-weight-modulated message passing (GCN, SAGE, GIN, APPNP, edge-weighted GAT; standard GAT and binary-mask architectures require modification); for contractive implicit models (IGNN-class), supplementary formal guarantees ($\ecrit$, three vulnerability regimes) additionally characterize convergence behavior. Predictions remain within 15\% of ground truth at $\varepsilon = 0.10$ across 7 architectures and 9 datasets, with continuous-to-discrete transfer positive in 88\% of architecture-dataset settings. The vulnerability spectrum correlates with power grid contingency rankings (P@10 = 0.66--0.87) as a proof-of-concept cross-domain demonstration. The per-edge ranking directly informs defense design (\cref{sec:defense_ablation}).

\textbf{Limitations.} (1)~First-order sensitivity thresholds $r_v$ are diagnostic quantities, not probabilistic certificates; they are locally tight for small $\varepsilon$ but can be violated at larger perturbation magnitudes (\cref{rem:certificates}). (2)~Continuous edge-weight perturbations; discrete insertions/deletions are outside the formal guarantee. Continuous-to-discrete transfer is positive in 29/33 architecture-dataset settings but negative in 4 (12\%), particularly for shallow models (GCN-2) and IGNN on dense graphs (\cref{tab:tau_cross}); practitioners should verify transfer for their specific setting. (3)~Standard GAT and architectures using binary adjacency masks require modification for continuous sensitivity; ``architecture-agnostic'' applies to GNNs with continuous edge-weight-modulated message passing. (4)~Power flow $\tau = 0.37$--$0.72$ is a proof-of-concept; operational deployment requires impedance-aware edge encoding, broader training data, and thermal/voltage violation metrics~\cite{wood2014power}. (5)~The matrix-free pipeline's Neumann series requires $\kappa < 1$; for explicit GNNs without contractivity, the dense path ($N \leq 200$) or finite-difference columns are used. (6)~IGNN accuracy (77.5\% on Cora) lags behind explicit GNNs ($\sim$82\% APPNP); this is the cost of contractivity constraints, not a limitation of $S_c$. The framework applies identically to higher-accuracy explicit models, which receive the full computational analysis but not the formal regime characterization. (7)~50-node subgraph rankings correlate weakly with full-graph rankings on large citation networks ($\tau = 0.16$ on Cora); the matrix-free full-graph pipeline is recommended for graph-level assessment.

\textbf{Ethical considerations and dual-use risks.} \AEGIS identifies maximally sensitive perturbation directions and structurally vulnerable edges, which could be misused to attack deployed GNNs in safety-critical domains---particularly concerning given the power grid case study, where attack optimization tools targeting critical infrastructure demand careful consideration. We argue the diagnostic value outweighs this risk: defenders must understand their system's attack surface to protect it, and the same vulnerability map directly informs defense design (\cref{sec:defense_ablation}). Responsible disclosure: we will release code with a tiered access protocol. Vulnerability analysis tools (per-edge rankings, sensitivity thresholds) will be openly available as they are purely defensive. The attack-generation component (SVD-identified perturbation direction) will require institutional affiliation for access, consistent with established practice in adversarial ML~\cite{wu2019adversarial}. Users of the power grid module should note that vulnerability rankings are conditioned on the GNN model's learned representation, not on physical grid parameters, and should not be used as a substitute for engineering contingency analysis.

\textbf{Future work.} Second-order bounds for formal certificates; discrete perturbation models that extend beyond the continuous threat model; integration with defense mechanisms~\cite{jin2020graph,zhang2020gnnguard} for vulnerability-aware training with dynamic retraining after edge protection; extending full-graph analysis beyond the current validated boundary ($N{=}7{,}650$ on a single GPU) via distributed JVP computation; edge-feature GNN architectures for power flow applications where line impedance, thermal rating, and voltage level modulate electrical coupling; and investigation of the connection between $S_c$ vulnerability scores and graph curvature-based bottleneck analysis~\cite{topping2022understanding}.
